\section{Related Work}
\label{sec:related}

\subsection{Video Diffusion Transformers}
\label{sec:rw_dits}

Diffusion models have become the dominant class of generative models for
high-resolution visual synthesis. Denoising diffusion probabilistic
models~\citep{ho2020ddpm} established the canonical training objective, and
latent diffusion~\citep{rombach2022ldm} brought the cost of high-resolution
synthesis into a practical range. Replacing the U-Net backbone with a
transformer, DiT~\citep{peebles2023dit} showed that diffusion models scale
predictably with width, depth, and tokens, and introduced AdaLN-Zero, a
per-block adaptive layer-norm modulation that conditions each transformer
block on a global timestep embedding.

The video setting introduces two pressures: high token counts from the
temporal dimension, and the need for long-horizon temporal coherence. Stable
Video Diffusion~\citep{blattmann2023svd} extends latent diffusion to
image-to-video with curated pretraining stages. CogVideoX~\citep{yang2024cogvideox}
and HunyuanVideo~\citep{kong2024hunyuanvideo} push the diffusion-transformer
formulation to 5B--13B parameters with full 3D attention and expert or
dual-stream layouts. Stable Diffusion 3 / MMDiT~\citep{esser2024sd3} couples
the DiT architecture with rectified-flow
training~\citep{lipman2023flowmatching}.
\backbone~\citep{meituan2025longcatvideo} is a 13.6B-parameter DiT trained
with rectified flow that supports text-to-video, image-to-video, and
video-continuation in a single set of weights, with minute-scale continuation
as part of its pretraining recipe.

Two architectural features recur across this family: a single global
timestep-embedding pathway, and per-block adaLN projections that map that
pathway into shift, scale, and gate vectors. Across these systems, the
modulation pathway is used purely at training time; adaptation after
pretraining is pursued through fine-tuning, LoRA, or full retraining rather
than through the modulation interface itself.

\subsection{Test-Time Adaptation and Test-Time Training}
\label{sec:rw_tta}

Test-time adaptation (TTA) and test-time training (TTT) form a closely
related family of methods that update a pretrained model on each test
instance or batch before producing predictions, without access to source
training data at adaptation time. TTT~\citep{sun2020ttt} constructs a
self-supervised auxiliary problem---e.g., rotation prediction---and takes a
short update on each unlabeled test sample before predicting.
Tent~\citep{wang2021tent} shows that, on classification under distribution
shift, fully test-time adaptation that updates only the channel-wise affine
parameters of normalization layers by minimizing prediction entropy is
sufficient to absorb large shifts. \citet{xiao2024ttasurvey} catalog more
than 400 follow-ups across model-update, inference-time, normalization,
sample-based, and prompt-based TTA variants.

A consistent finding in this literature is that very small, structured
adaptation surfaces---BatchNorm statistics, channel-wise affine parameters,
or a handful of feature-wise biases---often perform as well as or better
than freeing larger parts of the network. Almost all of this work, however,
is built on top of discriminative classifiers or segmentation models.
TTA-style methods that target large generative models---and generative
\emph{video} models in particular---remain comparatively rare. The few
entries in this intersection cluster around per-subject customization
(Section~\ref{sec:rw_personalization}) and around test-time \emph{search}
over noise trajectories~\citep{liu2025videoT1}, which is orthogonal to
updating either model weights or input embeddings.

\subsection{Adaptation and Personalization of Pretrained Diffusion Models}
\label{sec:rw_personalization}

A second body of work specializes a pretrained diffusion generator to
specific subjects, styles, or motions, typically with a small set of
reference images or a single reference video.
DreamBooth~\citep{ruiz2023dreambooth} fine-tunes the full text-to-image
diffusion network with a class-prior preservation loss to bind a new token
to a subject. Textual Inversion~\citep{gal2023textualinversion} takes the
opposite extreme, learning only a new text embedding while keeping the
diffusion network frozen. HyperDreamBooth~\citep{ruiz2024hyperdreambooth}
accelerates per-subject adaptation roughly $25\times$ by predicting an
initial low-rank weight update with a hypernetwork before a short
fast-finetune.

In the video setting, CustomTTT~\citep{bi2025customttt} combines per-LoRA
appearance and motion customization with an additional test-time training
pass to repair artifacts that appear when multiple LoRAs are composed.
``One-Minute Video Generation with Test-Time
Training''~\citep{dalal2025oneminuteTTT} inserts TTT layers into a
CogVideoX backbone so that hidden states themselves are small neural
networks, enabling long-context generation beyond what linear-attention
alternatives support. These methods share the broad goal of specializing a
frozen pretrained video diffusion transformer at inference time, but their
adaptation surfaces remain comparatively large---per-LoRA modules or entire
neural-network hidden states---and their target task is style or motion
transfer from a reference, rather than steering a continuation toward its
own conditioning frames.

\subsection{Parameter-Efficient Fine-Tuning}
\label{sec:rw_peft}

Parameter-efficient fine-tuning (PEFT) provides the standard family of
adaptation tools used throughout the diffusion ecosystem. Houlsby-style
adapters~\citep{houlsby2019adapters} insert small bottleneck modules
between frozen transformer sublayers. LoRA~\citep{hu2022lora} learns a
low-rank additive update to selected linear weights that can be merged
back into the base model at deployment. BitFit~\citep{zaken2022bitfit}
shows that updating only a model's bias terms---under $0.1\%$ of
parameters---can match full fine-tuning on language tasks.
SVDiff~\citep{han2023svdiff} fine-tunes only the singular values of each
weight matrix's SVD, yielding adaptation footprints on the order of a few
megabytes for text-to-image personalization.

Two empirical patterns recur when these methods are applied to large
diffusion backbones. First, the rank or dimensionality of the adaptation
subspace can be remarkably small while still producing useful adaptations.
Second, when the adaptation signal itself is small---a single clip rather
than a curated multi-image set---LoRA-style methods can either fail to
move the model in a useful direction or overfit within a small number of
gradient steps.

\subsection{Adaptive Normalization and Modulation Pathways}
\label{sec:rw_modulation}

Feature-wise Linear Modulation~\citep{perez2018film} introduced the
per-channel affine scale-and-shift conditioning operation that subsequent
adaptive-normalization variants inherit. Conditional batch and layer
normalization in diffusion models follow the same template, and AdaLN-Zero
in DiT~\citep{peebles2023dit} makes the timestep embedding the primary
modulator: each transformer block linearly projects a global conditioning
vector into shift, scale, and gate vectors for both attention and the MLP.
CogVideoX, HunyuanVideo, and \backbone retain this AdaLN-style modulation,
exposing a single \dimt-dimensional timestep embedding that is shared
across all transformer blocks. A small but growing line of work explores
making the modulation pathway itself more efficient or more controllable;
for example, AdaLN-LoRA, used inside Cosmos-style video DiTs, applies a
low-rank adapter to the adaLN projection during training. These works
modify the modulation pathway only at training time, and the pathway has
not, to our knowledge, been exposed as a per-instance adaptation interface
at inference.

\subsection{Evaluation of Video Generation}
\label{sec:rw_eval}

Fr\'echet Video Distance~\citep{unterthiner2018fvd} embeds videos through
an Inflated 3D ConvNet and computes a Fr\'echet distance between real and
generated activations; it remains the most widely reported single-number
video metric. Subsequent work shows that FVD is partly biased toward
per-frame content quality, and that swapping the I3D backbone for
self-supervised features mitigates this content
bias~\citep{ge2024contentdebiasedFVD}.
VBench~\citep{huang2024vbench} decomposes video quality into a battery of
fine-grained, human-aligned dimensions and is now the de facto benchmark
for text-to-video systems; VBench-2.0 extends the suite toward physics,
commonsense, and motion realism. Among datasets,
Panda-70M~\citep{chen2024panda70m} is a 70M-clip captioned video corpus
built from HD-VILA-100M with multiple cross-modality teacher captions, and
UCF-101~\citep{soomro2012ucf101} is a 101-class action-recognition dataset
that has been widely adopted as a continuation benchmark in prior
video-generation evaluation, where FVD is typically reported alongside
per-frame PSNR, SSIM, and LPIPS.
